
\def\PUDcodigo{ACE012}

\if\COMPILAR0
   \def\PUDcodacad{04507.18}
\else
   \def\PUDcodacad{04506.18}
\fi

\def\PUDdisciplina{Física 3}

\def\PUDsemestre{S4}

\def\PUDhoras{80}

%NOVOS ---------------------------------------------------
\def\PUDhorasTeoria{60}
\def\PUDhorasPratica{20}

\if\COMPILAR0
   \def\PUDinterECA{ 
   \hyperref[PUD:ECA210]{Circuitos I (04507.20)}
   
   \hyperref[PUD:ECA223]{Máquinas Elétricas (04507.36)} 
   }
\else
   \def\PUDinterECA{ 
   \hyperref[PUD:ECA210]{Circuitos I (04506.21)}
   
   \hyperref[PUD:ECA223]{Máquinas Elétricas (04506.36)} 
   }
\fi

\def\PUDatuaisECA{- Carregamento de baterias sem fio}

\def\PUDrecursos{Projetor multimídia; Quadro branco e pincel; Aulas em laboratório de eletroeletrônica.}

%----------------------------------------------------------

\def\PUDcreditos{4}

\def\PUDprerequisito{ \hyperref[PUD:ACE011]{ACE011} }

\if\COMPILAR0
   \def\PUDpreacad{ \hyperref[PUD:ACE011]{Física 2 (04507.13)} }
\else
   \def\PUDpreacad{ \hyperref[PUD:ACE011]{Física 2 (04506.13)} }
\fi

\def\PUDobjetivos{Compreender as etapas do método científico e estabelecer um diálogo com temas do cotidiano relacionados com fenômenos elétricos e magnéticos.}

\def\PUDementa{História da eletricidade e suas aplicações;  Carga elétrica e Campos elétricos;  Lei de Gauss;  Potencial Elétrico e Capacitância;  Corrente;  Campos Magnéticos e suas Fontes;  Indução e Indutância;  Análise das Oscilações eletromagnéticas;  Equações de Maxwell;  Magnetismo da Matéria.}

\def\PUDconteudo{ & História da eletricidade e suas aplicações: Textos. 
\\ 
 & Carga elétrica e Campos Elétricos: Condutores e não Condutores.  Lei de Coulomb e Formas de calcular o Campo Elétrico. 
\\ 
 & Lei de Gauss: Fluxo e Aplicações da Lei de Gauss. 
\\ 
 & Potencial Elétrico e Capacitância: Energia Potencial Elétrica.  Formas de Calcular o Potencial Elétrico.  Potencial de um Condutor Carregado. 
\\ 
 & Corrente, Resistores e Circuitos: Leis de Ohm.  Resistência.  Circuitos com uma e duas Malhas. 
\\ 
 & Campos Magnéticos e suas Fontes: Campo criado por um ímã e Corrente.  Lei de Ampére. \\ 
 & Indução e Indutância: Lei de Faraday.  Lei de Lenz.  Circuitos RL. 
\\ 
 & Análise das Oscilações eletromagnéticas: Oscilações em circuito LC.  Oscilações em circuito RLC.  Corrente alternada e transformadores. 
\\ 
 & Equações de Maxwell: Lei de Gauss para Campos Magnéticos.  Corrente de deslocamento. 
\\ 
 & Magnetismo da Matéria: Ímãs permanentes.  Diamagnetismo.  Paramagnetismo e ferromagnestismo. }

\def\PUDmetodo{Aulas expositórias que podem ser teóricas e/ou práticas, onde as práticas no laboratório serão marcadas no decorrer da disciplina.}

\def\PUDavalia{A avaliação será feita com aplicação de provas teóricas e/ou práticas, além da possibilidade de inclusão de trabalhos e seminários no decorrer da disciplina.}

\def\PUDbibbasica{ &  \href{www.biblioteca.ifce.edu.br/index.asp?codigo_sophia=23914}{Hugh}, D. Young; Roger, A. F.  Física 3 - Eletromagnetismo.  12º ed.  São Paulo, SP: Addison-Wesley, 2009.  425p.  ISBN: 9788588639348.
\\ 
 &  \href{www.biblioteca.ifce.edu.br/index.asp?codigo_sophia=41358}{Halliday}, D.; Resnick, R.; Walker, J.  Fundamentos de Física - Eletromagnetismo. v. 3.  9º ed.  Rio de Janeiro, RJ: LTC, 2013.  375p.  ISBN: 9788521619055.
\\ 
 & \href{www.biblioteca.ifce.edu.br/index.asp?codigo_sophia=65083}{Nussenzweig}, M. Curso de Física Básica 3: Eletromagnetismo.  2º ed.   São Paulo, SP: Edgard Blücher, 2015.  295p.  ISBN: 9788521208013. }

\def\PUDbibcomplementar{ & \href{www.biblioteca.ifce.edu.br/index.asp?codigo_sophia=63343}{Tipler},P.; Mosca, G.  Física para cientistas e engenheiros: eletricidade e magnetismo, óptica.  6º ed.  Vol 2.  Rio de Janeiro, RJ: LTC, 2015.  530p.  ISBN: 9788521617112.
\\ 
 & \href{www.biblioteca.ifce.edu.br/index.asp?codigo_sophia=24416}{Reitz}, O. R.; Milford, F. J.; Christy, R. W.  Fundamentos da Teoria Eletromagnética.  1º ed.  Rio de Janeiro, RJ: Elsevier: Campus, 1982.  516p.  ISBN: 8570011032.
\\ 
 & \href{www.biblioteca.ifce.edu.br/index.asp?codigo_sophia=62471}{Luiz}, A. M.  Física 3: eletromagnetismo, teoria e problemas resolvidos.  2º ed.  São Paulo, SP: Livraria da Física, 2009.  260p.  ISBN: 9788578610104.
\\
 & \href{www.biblioteca.ifce.edu.br/index.asp?codigo_sophia=42007}{Sadiku}, Matthew N. O.  Elementos de eletromagnetismo.  5º ed.  Porto Alegre, RS: Bookman, 2012.  702p.  ISBN: 9788540701502.
\\
 & \href{www.biblioteca.ifce.edu.br/index.asp?codigo_sophia=24465}{Wolski}, Belmiro.  Eletromagnetismo.  Curitiba, PR: Base Editorial, 2010.  128p.  ISBN: 9788579055515. }

\def\PUDautor{Francisco Frederico dos Santos Matos}

\def\PUDdata{2014-04-21}

\def\PUDrevisao{2}

\def\PUDrevisor{Celso Rogério Schmidlin Júnior}

\def\PUDdatarevisao{2019-05-15}

\PUDcorpo
